
% Stage23 frozen manuscript integration.
% Generated from sealed Stage23 JSON/CSV evidence only.

\subsection{Validation-Safe Shortcut-Feature Audit}

The full 70-feature ensemble achieved a random-natural PR-AUC of
0.995590 and ROC-AUC of 0.998625,
compared with chronological-natural PR-AUC 0.106215
and ROC-AUC 0.514918. The corresponding chronological
attack prevalence was 0.104847, placing the chronological PR-AUC only
0.001368 above its prevalence
reference.

Frozen primary ablations revealed metric-dependent split interactions. Some
pre-specified removals disproportionately reduced random-natural ranking, whereas
others affected chronological ranking more strongly. Matched-size placebo removals
also yielded non-zero interactions and bootstrap intervals excluding zero. Therefore,
interaction significance alone does not identify leakage.

The frozen depth-1 controls for Dst Port, Init Fwd Win Byts, and Fwd Seg Size Min
showed substantially stronger discrimination under random-natural validation than
under chronological-natural validation. These observations are consistent with
split-specific discriminative cues that transfer poorly, but they do not establish
causality.

Component-specific TreeSHAP analyses further demonstrated importance redistribution
after feature removal. Equal-weight normalized LightGBM/XGBoost shares are used only
as a descriptive consensus and are not exact SHAP values for the averaged-probability
ensemble.

The behavior-restricted representation retained strong random-natural performance but
did not preserve chronological discrimination. Attack-family conditioning further
showed that all 62,256 chronological attack rows belonged to
Infilteration, while random-natural validation contained 11 families meeting the frozen
interpretive support threshold.

Taken together, Stage23 supports a bounded conclusion of strong validation sensitivity
and poor forward-temporal transfer. It does not prove that any tested feature is leakage.
